% !TEX program = lualatex
\documentclass[aspectratio=43,professionalfonts,handout]{beamer}
\usepackage{beamer_template}

\newcommand{\LectureNum}{4}
\newcommand{\LectureTitle}{逆ラプラス変換}
\newcommand{\TextPages}{p.\,57-61}
\newcommand{\CourseName}{応用数学}
\renewcommand{\TermName}{前期}

\begin{document}

\begin{frame}{\CourseName\quad 第\LectureNum 回「\LectureTitle」}
  {\small \TermName\quad \TeacherName\quad ／\quad 教科書 \TextPages}

  \textbf{目標}：逆ラプラス変換の定義を理解し、部分分数分解を用いて逆変換を求めることができる

\end{frame}

\begin{frame}{逆ラプラス変換の定義}
  \begin{block}{}
  \begin{itemize}
    \item $\mathcal{L}^{-1}[F(s)] = f(t)$
    \item 原関数の一致定理
    \item 逆変換の線形性
  \end{itemize}
  \end{block}
\end{frame}

\begin{frame}{基本公式}
  \begin{block}{}
  \begin{itemize}
    \item $\mathcal{L}^{-1}[1/(s-\alpha)] = e^{\alpha t}$
    \item $\mathcal{L}^{-1}[1/(s-\alpha)^{2}] = te^{\alpha t}$
    \item $\mathcal{L}^{-1}[1/((s-\alpha)(s-\beta))] = (e^{\alpha t}-e^{\beta t})/(\alpha-\beta)$
    \item $\mathcal{L}^{-1}[1/((s-\alpha)^{2}+\beta^{2})] = (1/\beta)e^{\alpha t}\sin \beta t$
  \end{itemize}
  \end{block}
\end{frame}

\begin{frame}{部分分数分解}
  \begin{block}{}
  \begin{itemize}
    \item 手順：分母を因数分解  $\to$  分解  $\to$  係数決定  $\to$  逆変換
    \item 例題14(1)：s/((s+2)(s+3)(s+4))
    \item 例題14(2)：$(3s^{2}+5)/((s+1)^{2}(s-3))$
  \end{itemize}
  \end{block}
\end{frame}

\begin{frame}{演習（練習問題1）}
  \begin{exampleblock}{自分で解いてみよう}
  \begin{itemize}
    \item p.61 問1〜5より選題
  \end{itemize}
  \end{exampleblock}
\end{frame}

\begin{frame}{まとめ・前半総括}
  \begin{itemize}
    \item 逆ラプラス変換の手順
    \item ラプラス変換表の活用
    \item 次回：微分方程式への応用
  \end{itemize}
\end{frame}

\end{document}
